\documentclass[11pt,a4paper]{article}
\usepackage[margin=2.3cm]{geometry}
\usepackage{amsmath,amssymb,bm}
\usepackage{booktabs}
\usepackage{array}
\usepackage{xcolor}
\usepackage{hyperref}
\usepackage{enumitem}
\usepackage[most]{tcolorbox}
\usepackage{longtable}
\usepackage{microtype}

\hypersetup{
    colorlinks=true,
    linkcolor=blue!60!black,
    urlcolor=blue!60!black,
    citecolor=blue!60!black
}

\definecolor{deepblue}{RGB}{24,64,120}
\definecolor{softblue}{RGB}{235,245,255}
\definecolor{softorange}{RGB}{255,246,235}
\definecolor{softpurple}{RGB}{245,240,255}

\setlist[itemize]{leftmargin=2em}
\setlist[enumerate]{leftmargin=2em}

\title{A Unified Framework for Nonlinear Local--Global Latent Dynamics,\\ Koopman Modal Signatures, and Two LFP--BOLD Projects}
\author{Framework note}
\date{2026-06-10}

\begin{document}
\maketitle

\begin{tcolorbox}[colback=softblue,colframe=deepblue,title=Core idea]
LFP, local BOLD, and global BOLD should not be treated as a simple feed-forward input-output chain. They are better understood as different observable readouts of a nonlinear local-global latent system. Koopman eigenfunctions transform this complex nonlinear system into interpretable modal variables, making the global gating state, pre-activation modes, global ROI maps, and loop backbones operational phenomenological modal signatures.
\end{tcolorbox}

\tableofcontents
\newpage

\section{Overview}

This note summarizes the current mathematical framework connecting two related projects:
\begin{enumerate}
    \item \textbf{Event-related LFP subprocess $\to$ global BOLD pattern / BOLD eigenfunction residual}: how local LFP subprocesses can act as intrinsic perturbations that trigger or predict global BOLD modal innovations.
    \item \textbf{Non-causal LFP--BOLD impulse responses / pre-activation}: why local BOLD can show apparent pre-activation relative to local LFP, and how this can be interpreted using closed-loop local-global latent dynamics and a Koopman modal representation.
\end{enumerate}

The two projects correspond to two halves of the same closed loop:
\[
\underbrace{u_{\mathrm{LFP}} \rightarrow x^{(g)}}_{\text{local-to-global: BOLD modal innovation}}
\qquad\text{and}\qquad
\underbrace{x^{(g)} \rightarrow u_{\mathrm{LFP}}}_{\text{global-to-local: pre-activation / gating}}.
\]
Together, they describe
\[
x^{(g)} \rightarrow u_{\mathrm{LFP}} \rightarrow x^{(g)},
\]
meaning that a global order-parameter-like state gates local hippocampal subprocesses, and local subprocesses in turn perturb or update the global BOLD network state.

\section{Underlying model: a nonlinear local-global latent system}

Let the latent state be decomposed into local and global components:
\begin{equation}
    x_t =
    \begin{bmatrix}
    x_t^{(l)} \\
    x_t^{(g)}
    \end{bmatrix}.
\end{equation}

Here:
\begin{itemize}
    \item $x_t^{(l)}$ is a \textbf{local latent state}, such as hippocampal local neural state, local vascular state, or local circuit excitability.
    \item $x_t^{(g)}$ is a \textbf{global latent state / internal environment / global order-parameter-like state}, such as global brain state, neuromodulatory state, brainstem-thalamic-cortical state, global vascular state, or arousal state.
    \item $u_t$ is a \textbf{local LFP subprocess / local neural action}.
\end{itemize}

A generic nonlinear model can be written as
\begin{align}
    x_{t+1}^{(l)} &= F_l\!\left(x_t^{(l)}, x_t^{(g)}, u_t\right), \\
    x_{t+1}^{(g)} &= F_g\!\left(x_t^{(g)}, x_t^{(l)}, u_t\right),
\end{align}
or more compactly:
\begin{equation}
    x_{t+1}=F(x_t,u_t).
\end{equation}

The key point is not that LFP is an external stimulus. Rather, $u_t$ is treated as an action-like subprocess of the local subsystem. It can be written as an input-like variable for modeling convenience, but in the closed-loop view it is generated internally by the system.

\section{BOLD is a nonlinear readout, not the mechanism itself}

Local BOLD and global BOLD are observable readouts of the latent state:
\begin{align}
    y_t^{(lB)} &= h_l\!\left(x_t^{(l)}, x_t^{(g)}\right), \\
    y_t^{(gB)} &= h_g\!\left(x_t^{(l)}, x_t^{(g)}\right).
\end{align}

Therefore, local BOLD is not simply the output of a local HRF driven by local LFP. It may contain both a local and a global component:
\begin{equation}
    y_t^{(lB)} = \text{local component} + \text{global component}.
\end{equation}

This is crucial: if local BOLD reads out a global latent state before that state gates a future local LFP subprocess, local BOLD can show apparent pre-activation relative to LFP.

\section{Global-to-local path: global state gates local LFP subprocesses}

The central explanation for pre-activation is:
\begin{equation}
    x_t^{(g)} \rightarrow u_{\mathrm{LFP},t+\Delta}.
\end{equation}

More explicitly:
\begin{equation}
    x_t^{(g)} \rightarrow x_{t+\Delta}^{(l)} \rightarrow u_{\mathrm{LFP},t+\Delta}.
\end{equation}

This means that a global latent state / order-parameter-like variable appears first, then modulates the local hippocampal condition, making a local LFP subprocess more likely to occur.

To form a complete closed loop, there may also be a local-to-global update:
\begin{equation}
    u_{\mathrm{LFP},t} \rightarrow x_{t+\Delta}^{(g)}.
\end{equation}

The overall loop can therefore be summarized as
\begin{equation}
    x^{(g)} \rightarrow u_{\mathrm{LFP}} \rightarrow x^{(g)}.
\end{equation}

However, for explaining \textbf{pre-activation}, the most important part is the first half:
\begin{equation}
    x^{(g)} \rightarrow u_{\mathrm{LFP}}.
\end{equation}

\section{Why local BOLD--LFP can produce an apparent non-causal IR}

The standard feed-forward HRF model assumes:
\begin{equation}
    y_{\mathrm{localBOLD}}(t)
    =
    \sum_{\tau\ge 0} h(\tau) u_{\mathrm{LFP}}(t-\tau)+\epsilon(t).
\end{equation}

This treats LFP as an exogenous input and BOLD as a passive output. Under this assumption, if the estimated impulse response has negative-lag pre-activation, it appears non-causal.

In the local-global latent model, a global state $x_t^{(g)}$ can simultaneously:
\begin{enumerate}
    \item be read out early by local BOLD:
    \begin{equation}
        x_t^{(g)} \rightarrow y_t^{(lB)},
    \end{equation}
    \item gate a future local LFP subprocess:
    \begin{equation}
        x_t^{(g)} \rightarrow u_{\mathrm{LFP},t+\Delta}.
    \end{equation}
\end{enumerate}

Thus the observations may show
\begin{equation}
    y_{\mathrm{localBOLD},t}
    \quad \text{preceding} \quad
    u_{\mathrm{LFP},t+\Delta}.
\end{equation}

If one still applies a feed-forward LFP-to-BOLD deconvolution model, the resulting cross-kernel can show negative-lag pre-activation. This does \textbf{not} imply
\begin{equation}
    y_{\mathrm{localBOLD}} \rightarrow u_{\mathrm{LFP}}.
\end{equation}
Instead, it reflects the shared influence of a latent global state:
\begin{equation}
    x^{(g)} \rightarrow
    \begin{cases}
    y_{\mathrm{localBOLD}},\\
    u_{\mathrm{LFP}}.
    \end{cases}
\end{equation}

The apparent non-causality therefore arises when a closed-loop endogenous system is interpreted using an open-loop feed-forward model.

\section{From nonlinear system to augmented autonomous system}

If $u_t$ is generated internally, it can be written as
\begin{equation}
    u_t = \pi\!\left(x_t^{(g)}, x_t^{(l)}, u_{t-1},u_{t-2},\ldots\right)+\epsilon_t.
\end{equation}

By augmenting the state with $u_t$ and its history,
\begin{equation}
    s_t = \begin{bmatrix}
    x_t \\
    u_t \\
    u_{t-1} \\
    \vdots
    \end{bmatrix},
\end{equation}
the closed-loop system can be written in autonomous form:
\begin{equation}
    s_{t+1}=F_{\mathrm{cl}}(s_t).
\end{equation}

The readouts can be written as
\begin{align}
    u_t &= C_u s_t, \\
    y_t^{(lB)} &= C_l s_t, \\
    y_t^{(gB)} &= C_g s_t.
\end{align}

This step is important because it shows that the local LFP subprocess, local BOLD, and global BOLD can all be treated as different readouts of the same autonomous closed-loop system.

\section{Koopman modal representation}

For the augmented autonomous system
\begin{equation}
    s_{t+1}=F_{\mathrm{cl}}(s_t),
\end{equation}
let $a_t$ denote Koopman eigenfunction coordinates. In modal coordinates:
\begin{equation}
    a_{t+1}=\Lambda a_t.
\end{equation}

The observables can be approximated as linear readouts:
\begin{align}
    u_{\mathrm{LFP},t}&=C_u a_t, \\
    y_{\mathrm{localBOLD},t}&=C_l a_t, \\
    y_{\mathrm{globalBOLD},t}&=C_g a_t.
\end{align}

The value of the Koopman representation is that the original latent dynamics and observation functions may be nonlinear, but in lifted/modal space the dominant modes, phase relationships, and spatial readouts can be analyzed linearly or approximately linearly.

\section{Definition of pre-activation modes}

For an LFP--BOLD anchor pair $(p,q)$:
\begin{itemize}
    \item $p$ denotes a local LFP band/subprocess, such as PGO, theta, or ripple.
    \item $q$ denotes a local BOLD ROI or local BOLD readout.
\end{itemize}

If the empirical cross-kernel $h_{q,p}(\tau)$ shows pre-activation at $\tau<0$, we can identify Koopman modes that contribute to this negative-lag component.

Define:
\begin{equation}
\mathcal I_{\mathrm{pre}}(q,p)
=
\{i: \text{mode } i \text{ contributes to the negative-lag component of pair }(q,p)\}.
\end{equation}

These are the \textbf{pre-activation modes} for the anchor pair.

Biological interpretation:
\begin{quote}
These modes most naturally correspond to a global latent gating state / order-parameter-like variable. This global state is read out by BOLD before it gates future local LFP subprocesses.
\end{quote}

Important boundary:
\begin{quote}
Pre-activation modes are not causal modes. They are pair-conditioned modal signatures.
\end{quote}

\section{Koopman phase-lead explanation}

For mode $i$, let the continuous-time eigenvalue be
\begin{equation}
    \nu_i = \alpha_i + \mathrm{i}\omega_i.
\end{equation}

The complex readout weights for LFP and local BOLD can be written as
\begin{align}
    C_u(i)&=|C_u(i)|e^{\mathrm{i}\phi_{u,i}}, \\
    C_l(i)&=|C_l(i)|e^{\mathrm{i}\phi_{l,i}}.
\end{align}

If the BOLD readout phase leads the LFP readout phase, this mode contributes to apparent pre-activation in the LFP--BOLD cross-kernel. A phase-derived time difference can be written as
\begin{equation}
    \Delta t_{lB-u,i}
    =
    \frac{\operatorname{wrap}(\phi_{l,i}-\phi_{u,i})}{\omega_i},
\end{equation}
with the exact sign depending on the phase convention. Conceptually:
\begin{equation}
\text{BOLD phase leads LFP phase}
\quad \Rightarrow \quad
\text{negative-lag contribution in empirical IR}.
\end{equation}

\section{Meaning of the global ROI map}

For the pre-activation modes of an anchor pair $(q,p)$, define the global ROI map as
\begin{equation}
M_{q,p}(r)
=
\sum_{i\in\mathcal I_{\mathrm{pre}}(q,p)} s_i |C_g(r,i)|,
\end{equation}
where:
\begin{itemize}
    \item $r$ is a global BOLD ROI.
    \item $C_g(r,i)$ is the global BOLD readout weight of mode $i$ at ROI $r$.
    \item $s_i$ is the strength with which mode $i$ contributes to the pair-specific pre-activation, for example negative-lag area, mode amplitude, or pairwise mode weight.
\end{itemize}

Interpretation:
\begin{quote}
The global ROI map is the spatial footprint of the pre-activation modes in whole-brain BOLD ROI space.
\end{quote}

In the closed-loop / slaving-principle interpretation, it further corresponds to:
\begin{quote}
the BOLD-observable footprint of the global latent state / order-parameter-like variable that gates the local LFP subprocess.
\end{quote}

It is \textbf{not} a causal map of BOLD ROIs driving LFP. It is a map of which ROIs jointly express the global gating state.

\section{Meaning of the loop backbone}

For the same pre-activation modes, each ROI readout has both amplitude and phase:
\begin{equation}
    C_g(r,i)=|C_g(r,i)|e^{\mathrm{i}\phi_{r,i}}.
\end{equation}

The global ROI map uses the amplitude:
\begin{equation}
    |C_g(r,i)|,
\end{equation}
whereas the loop backbone uses the phase:
\begin{equation}
    \phi_{r,i}.
\end{equation}

Thus:
\begin{quote}
The global ROI map describes where the global state is expressed. The loop backbone describes how that global state is phase-expressed across ROIs.
\end{quote}

For each mode, one may define a phase-based adjacency:
\begin{equation}
A_i(r_1,r_2)=1
\quad \text{if ROI } r_1 \text{ phase-leads ROI } r_2
\quad \text{within mode } i.
\end{equation}

Aggregating over pre-activation modes gives
\begin{equation}
A_{q,p}^{\mathrm{backbone}}(r_1,r_2)
=
\sum_{i\in\mathcal I_{\mathrm{pre}}(q,p)} s_i A_i(r_1,r_2).
\end{equation}

The loop backbone is interpreted as:
\begin{quote}
the phase-ordered expression, or phase portrait, of the global gating state in ROI space.
\end{quote}

It is a candidate biological representation / phenomenological modal signature. It is not the anatomical feedback loop itself, and it does not by itself establish interventional causality.

\section{Relation to Haken's slaving principle}

This framework can be related to Haken's slaving principle:
\begin{equation}
\text{distributed brain regions}
\Rightarrow
\text{global order parameter}
\Rightarrow
\text{local fast subprocess}.
\end{equation}

In the current setting:
\begin{equation}
\text{distributed ROI-level activity}
\Rightarrow
x^{(g)}
\Rightarrow
u_{\mathrm{LFP}}.
\end{equation}

Here:
\begin{itemize}
    \item $x^{(g)}$ is a global order-parameter-like state / slaving variable.
    \item The global ROI map is the spatial BOLD footprint of $x^{(g)}$.
    \item The loop backbone is the phase portrait of $x^{(g)}$ across ROIs.
    \item The local LFP subprocess is the fast local process gated or enslaved by $x^{(g)}$.
\end{itemize}

This also forms a kind of circular causality:
\begin{equation}
\text{distributed global state}
\rightarrow
\text{local LFP subprocess}
\rightarrow
\text{global state update}.
\end{equation}

\section{Relationship to the two projects}

\subsection{Project 1: Event-related LFP subprocess $\to$ global BOLD pattern / BOLD eigenfunction residual}

This project corresponds to the local-to-global path:
\begin{equation}
    u_{\mathrm{LFP}} \rightarrow x^{(g)} \rightarrow y_{\mathrm{globalBOLD}}.
\end{equation}

Operationally, one first builds a BOLD-only Koopman model:
\begin{equation}
    b_{k+1}=\Lambda_B b_k + \eta_k^B,
\end{equation}
where the BOLD modal innovation / BOLD eigenfunction residual is
\begin{equation}
    \eta_k^B=b_{k+1}-\Lambda_B b_k.
\end{equation}

The high-rate LFP subprocesses are then converted to the BOLD time scale as density variables:
\begin{equation}
D_r^{\mathrm{LFP}}[k]
=
\text{density / occupancy of LFP subprocess } r
\text{ in BOLD window } k.
\end{equation}

The core predictive model is
\begin{equation}
\eta_{m,k}^B
=
\sum_r\sum_{\ell}
\beta_{m,r,\ell}D_r^{\mathrm{LFP}}[k-\ell]
+
\epsilon_{m,k}.
\end{equation}

If $D_r^{\mathrm{LFP}}$ predicts future $\eta^B$, then:
\begin{quote}
the local LFP subprocess acts as an intrinsic perturbation that triggers or predicts global BOLD modal innovations.
\end{quote}

Important boundary:
\begin{quote}
The BOLD residual itself is not the LFP subprocess activation. It is an input-like deviation of BOLD modal dynamics. It supports the LFP-subprocess-as-intrinsic-perturbation interpretation only if it is predicted by LFP subprocess density.
\end{quote}

\subsection{Project 2: Non-causal LFP--BOLD IR / pre-activation}

This project corresponds to the global-to-local path:
\begin{equation}
    x^{(g)} \rightarrow u_{\mathrm{LFP}}.
\end{equation}

The proposed sequence is:
\begin{enumerate}
    \item A global latent state $x^{(g)}$ appears first.
    \item Local BOLD reads out a local/global mixture that includes $x^{(g)}$.
    \item The same $x^{(g)}$ gates a future local LFP subprocess.
    \item Therefore a feed-forward LFP-to-local-BOLD IR can show apparent pre-activation.
\end{enumerate}

In the Koopman analysis:
\begin{itemize}
    \item The anchor pair $(q,p)$ selects the pre-activation modes $\mathcal I_{\mathrm{pre}}(q,p)$.
    \item The global ROI map $M_{q,p}(r)$ is the spatial footprint of these modes.
    \item The loop backbone $A_{q,p}^{\mathrm{backbone}}$ is their phase-ordered expression across ROIs.
\end{itemize}

This project argues:
\begin{quote}
BOLD pre-activation is not reverse neurovascular causality. It is a phenomenological modal signature of a global gating state that precedes and organizes local LFP subprocesses.
\end{quote}

\section{Integrated closed-loop interpretation}

The two projects correspond to two halves of the same closed loop:

\begin{longtable}{p{0.28\textwidth}p{0.32\textwidth}p{0.32\textwidth}}
\toprule
\textbf{Path} & \textbf{Project} & \textbf{Interpretation} \\
\midrule
$u_{\mathrm{LFP}} \rightarrow x^{(g)}$ & Event-related LFP subprocess $\to$ BOLD innovation & Local subprocess perturbs / updates global BOLD modal state \\
$x^{(g)} \rightarrow u_{\mathrm{LFP}}$ & Non-causal IR / pre-activation & Global gating state precedes and gates local LFP subprocess \\
\bottomrule
\end{longtable}

Together:
\begin{equation}
    x^{(g)} \rightarrow u_{\mathrm{LFP}} \rightarrow x^{(g)}.
\end{equation}

In words:
\begin{quote}
A global order-parameter-like state gates local hippocampal subprocesses, and local subprocesses in turn perturb or update the global BOLD network state.
\end{quote}

This closed-loop local-global self-organization is the shared mechanistic interpretation of the two projects.

\section{Interpretation boundaries}

The framework should be stated carefully:
\begin{itemize}
    \item A \textbf{pre-activation mode} is not a causal mode; it is a pair-conditioned modal signature.
    \item A \textbf{global ROI map} is not a causal feedback map; it is the spatial footprint of pre-activation modes.
    \item A \textbf{loop backbone} is not an anatomical feedback loop; it is the phase-ordered expression of a global gating state in ROI space.
    \item The whole interpretation is phenomenological but mechanistically informative.
\end{itemize}

A concise manuscript-friendly formulation is:
\begin{quote}
Pre-activation modes identify a global order-parameter-like state associated with an LFP--BOLD anchor pair. The global ROI map describes where this state is expressed in BOLD space, while the loop backbone describes how that state is phase-expressed across ROIs. These objects are phenomenological modal signatures of a latent feedback/slaving process, rather than direct anatomical or interventional feedback pathways.
\end{quote}

\end{document}
